% =============================================================
% Capitulo 2 - Fundamentos de la tecnologia LiDAR y HSRL
% =============================================================
\chapter{Fundamentos de la tecnolog\'ia LiDAR y HSRL}
\label{cap:lidar_hsrl}

\section{Principio de funcionamiento de un LiDAR}
\label{sec:principio_lidar}

Un LiDAR (\emph{Light Detection and Ranging}) es un instrumento de
teledetecci\'on activa que emplea radiaci\'on \'optica como se\~nal de
sondeo, de manera an\'aloga a como el radar emplea ondas de radio
\cite{weitkamp2005}. En su configuraci\'on b\'asica, un sistema LiDAR
consta de un transmisor -- t\'ipicamente un l\'aser pulsado, en
ocasiones seguido de un expansor de haz -- y un receptor, compuesto por
un telescopio que colecta los fotones retrodispersados por la
atm\'osfera, un sistema \'optico de an\'alisis y un detector que
convierte la se\~nal \'optica recibida en una se\~nal el\'ectrica
\cite{weitkamp2005}.

La potencia recibida por el sistema puede expresarse, en su forma m\'as
simple, como

\begin{equation}
  P(R) = K \, G(R) \, \beta(R) \, T(R)
  \label{eq:lidar_simple}
\end{equation}

\noindent donde $K$ resume el desempe\~no del sistema, $G(R)$ describe
la geometr\'ia de medici\'on en funci\'on de la distancia $R$, $\beta(R)$
es el coeficiente de retrodispersi\'on atmosf\'erico y $T(R)$ el
t\'ermino de transmisi\'on, que describe la atenuaci\'on de la se\~nal
en su trayecto de ida y vuelta \cite{weitkamp2005}. Una forma m\'as
detallada de la ecuaci\'on del LiDAR, que expl\'icita el \'area del
receptor $A$, la eficiencia del receptor $\eta$, la duraci\'on del
pulso $\tau$ y la funci\'on de solapamiento $O(r)$, es

\begin{equation}
  P(r) = P_0 \, \eta \, \frac{A}{r^2} \, \frac{c \tau}{2} \, O(r) \,
  \beta(r) \, \exp\left(-2\int_0^{r}\alpha(r')\,dr'\right)
  \label{eq:lidar_equation}
\end{equation}

\noindent siendo $P_0$ la potencia media transmitida durante el pulso,
$c$ la velocidad de la luz y $\alpha(r)$ el coeficiente de extinci\'on
atmosf\'erico \cite{weitkamp2005,kovalev2015}.

% PLACEHOLDER: si se requiere, ampliar con el desarrollo historico del
% LiDAR y el detalle de tecnicas (elastico, Raman, DIAL, Doppler) a
% partir del Capitulo 1 de \cite{weitkamp2005}.

\section{LiDAR de Alta Resoluci\'on Espectral (HSRL)}
\label{sec:hsrl}

Los coeficientes $\beta$ y $\alpha$ de la Ecuaci\'on~\ref{eq:lidar_equation}
se descomponen, cada uno, en una componente molecular y una componente
de aerosol:

\begin{align}
  \beta &= \beta_{mol} + \beta_{aer} \\
  \alpha &= \alpha_{mol,sca} + \alpha_{mol,abs} + \alpha_{aer,sca} + \alpha_{aer,abs}
\end{align}

\noindent lo que deja, en principio, seis inc\'ognitas a resolver a
partir de una \'unica ecuaci\'on de se\~nal \cite{weitkamp2005}. Un
LiDAR de Alta Resoluci\'on Espectral (HSRL) resuelve este problema
explotando el corrimiento Doppler que experimenta la luz dispersada por
mol\'eculas en movimiento t\'ermico aleatorio: la distribuci\'on
Maxwelliana de velocidades moleculares produce corrimientos del orden
de 1~GHz, mientras que los aerosoles y part\'iculas, cuyo movimiento
est\'a determinado por el viento y la turbulencia, producen
corrimientos mucho menores ($\sim$30~MHz y $\sim$3~MHz, respectivamente)
\cite{weitkamp2005}. El
espectro de la luz retrodispersada resulta as\'i de la superposici\'on
de un pico angosto debido a la dispersi\'on de part\'iculas sobre una
distribuci\'on mucho m\'as ancha producida por la dispersi\'on
molecular \cite{weitkamp2005}.

Utilizando un filtro \'optico angosto centrado en la longitud de onda
del l\'aser transmisor, es posible separar la contribuci\'on molecular
de la contribuci\'on de aerosol, obteniendo dos ecuaciones de LiDAR
independientes cuyo cociente permite determinar la raz\'on de
retrodispersi\'on $\Re(r) = \beta_{aer}(r)/\beta_{mol}(r)$ sin necesidad
de recurrir a hip\'otesis adicionales sobre la relaci\'on entre
extinci\'on y retrodispersi\'on de aerosoles, a diferencia de otros
m\'etodos de inversi\'on como el de Klett \cite{weitkamp2005,kovalev2015}.
Esta condici\'on de filtrado exige que el ancho de banda del filtro sea
del orden de 1~GHz y que la frecuencia del l\'aser
transmisor permanezca enclavada (\emph{locked}) a la frecuencia central
del filtro, con un ancho de l\'inea menor al del propio filtro
\cite{weitkamp2005}.

\section{Absorci\'on molecular de iodo como filtro de referencia}
\label{sec:absorcion_iodo}

Entre las implementaciones de HSRL existentes, el uso de filtros de
absorci\'on at\'omica o molecular constituye una alternativa a los
interfer\'ometros de Fabry-P\'erot, evitando la sensibilidad
t\'ermica y mec\'anica de estos \'ultimos \cite{weitkamp2005}. En
particular, el reemplazo de las celdas de vapor de bario por celdas de
iodo molecular permite operar a temperaturas sustancialmente menores
(entre $\sim$25\textcelsius\ y $\sim$100\textcelsius), conservando
la estabilidad espectral y el amplio \'angulo de aceptaci\'on de las
celdas at\'omicas, adem\'as de contar con l\'ineas de absorci\'on
convenientes dentro del rango de sintonizaci\'on t\'ermica de l\'aseres
Nd:YAG duplicados en frecuencia (532~nm)
\cite{weitkamp2005}.

En el sistema HSRL de Pilar, la celda de iodo cumple la funci\'on de
filtro rechaza-banda \'optico sintonizado naturalmente a una frecuencia
determinada. Las condiciones de operaci\'on del lazo son las siguientes
\cite{galvagno2026ilrc}:

\begin{itemize}
  \item La temperatura de la celda se mantiene entre 30\textcelsius\
        y 50\textcelsius\ para mantener el iodo en fase gaseosa.
  \item Temperaturas m\'as elevadas incrementan la concentraci\'on de
        iodo activo, lo que resulta en l\'ineas de absorci\'on m\'as
        intensas.
  \item Dos moduladores acusto\'opticos (AOM) generan los haces $P_+$ y
        $P_-$, desplazados en frecuencia $\pm$260~MHz respecto del haz
        principal $P_0$.
  \item Un l\'aser seeder sintoniza la longitud de onda del l\'aser
        host a una de las l\'ineas de absorci\'on de iodo m\'as
        intensas.
  \item En el pico de absorci\'on, los haces $P_+$ y $P_-$ presentan
        intensidades iguales, produciendo amplitudes id\'enticas en los
        fotodetectores.
\end{itemize}

\section{Sistema HSRL del Observatorio Geof\'isico y Meteorol\'ogico de Pilar}
\label{sec:hsrl_pilar}

El sistema instalado en Pilar sintoniza en longitud de onda un l\'aser
de granate de itrio y aluminio dopado con neodimio (Nd:YAG), denominado
``host'', mediante un l\'aser ``seeder'' de menor potencia acoplado por
fibra \'optica; la longitud de onda sintonizada depende
proporcionalmente de la temperatura de cavidad del seeder
\cite{sat2026}. En el lazo de control, un separador de haz toma una
muestra del l\'aser host, que luego se divide en dos haces id\'enticos
($P_0$) mediante otro separador. Estos dos haces son modulados en
amplitud por los AOM mencionados en la Secci\'on~\ref{sec:absorcion_iodo},
generando las se\~nales $P_+$ y $P_-$ que atraviesan la celda de iodo
\cite{sat2026}. El sistema encuentra el punto de trabajo correcto
cuando la temperatura del seeder es tal que la amplitud de la se\~nal
$P_0$ resulta m\'inima, es decir, cuando las se\~nales $P_+$ y $P_-$ se
igualan \cite{sat2026}.

Al finalizar la PPS, este lazo de control se operaba mediante un
osciloscopio y una PC \cite{pps2025}; el presente proyecto reemplaza
dicho conjunto por la placa de control microcontrolada descrita en la
Parte~II de este informe.
